\documentclass[10pt,a4paper]{amsart}
\usepackage[margin=19mm]{geometry}
\usepackage{amsmath,amssymb,mathtools}
\usepackage[scr=rsfs,cal=euler]{mathalfa}
\usepackage{xfrac}
\usepackage[unicode,hidelinks,bookmarks=false]{hyperref}
\newcommand{\ie}{i.e.\ }
\newcommand{\cf}{cf.\ }
\newcommand{\resp}{respectively\ }
\newcommand{\Subsection}{\subsection}
\providecommand{\nobreakdash}{\nobreak\hbox{-}}
\setlength{\emergencystretch}{2em}
\multlinegap0pt
\usepackage{bm}
\usepackage{bbm}
\usepackage{dsfont}
\usepackage{xspace}
\usepackage{cancel}
\usepackage{mathtools}
\usepackage{amsthm}
\makeatletter
\@ifundefined{theorem}{\newtheorem{theorem}{Theorem}}{}
\@ifundefined{lemma}{\newtheorem{lemma}[theorem]{Lemma}}{}
\makeatother
\def \R{\mathbb{R}}
\def\co{\colon} 
\title[A homotopy invariant of stable maps to oriented surfaces]{A homotopy invariant of stable maps to oriented surfaces}
\author{Liam Kahmeyer, Rustam Sadykov}
\date{Source: arXiv:2208.07297v3 (CC BY 4.0)}
\begin{document}
\maketitle

\noindent\emph{Source and licence.} A homotopy invariant of stable maps to oriented surfaces. Version arXiv:2208.07297v3 (CC BY 4.0), distributed under the Creative Commons Attribution 4.0 International licence, as recorded in the arXiv licence metadata for this exact version and shown on the abstract page https://arxiv.org/abs/2208.07297v3. Full rights notice: licenses/arxiv-2208.07297-CC-BY-4.0.txt.\par\smallskip

\noindent\textit{Editorial scope.} Counted unit: the Theorem whose source label is \texttt{th:4} in the article (source0.tex lines 453--454, source label \texttt{th:4}), delivered together with the article's own complete original proof of it (source0.tex lines 455--495, one \texttt{proof} environment of 41 lines). No mathematics is written, repaired or re-derived: statement and proof are the article's own text line for line. Required context retained uncounted: none. Every definition, display and label of those lines is retained unchanged. Omitted from this package and not redistributed: 1--452, 496--1380. Nothing inside a retained excerpt is omitted or altered: every retained body is delivered exactly as the article's lines are written, so no omission receipt of any kind accompanies this package. Citations: the retained excerpts carry no citation command at all, so no bibliography entry of the article is transcribed and no reference is redistributed. Reference adaptation: none was needed and none was made; every reference inside the delivered unit resolves inside it, so no unresolved reference or citation is produced. Printed slips kept literal: no printed word, symbol, hypothesis or number of the counted statement or of its proof was altered. Layout and class: the article's own class is \texttt{amsart} with packages including ulem, tikz-cd, graphicx, mathabx, geometry, amsthm; the wrapper is the lane's pinned \texttt{amsart} editorial wrapper at 10pt on A4, which restates the theorem-like environments the retained text needs and adds the article's own macro definitions that the retained text needs (\texttt{\textbackslash R}, \texttt{\textbackslash co}), with extra wrapper packages (bm, bbm, dsfont, xspace, cancel, mathtools). The delivered numbering is the unit's own. Assets: no retained span contains a figure environment, a graphics-inclusion command, a PDF-page inclusion, a table or a TikZ picture; no image file is copied into this package, the package declares no asset dependency and no image file is redistributed with it. Licence: version arXiv:2208.07297v3 of this article is distributed under the Creative Commons Attribution 4.0 International licence, as recorded in the arXiv licence metadata for this exact version and shown on the abstract page https://arxiv.org/abs/2208.07297v3. A full rights notice is delivered at licenses/arxiv-2208.07297-CC-BY-4.0.txt. Characters: every retained excerpt is pure ASCII, so the delivered engine, XeLaTeX with its default Latin Modern text font, reports no missing character.
\begin{theorem}\label{th:4}  Under a generic homotopy $F = \{f_t\}$ of maps to $\R^2$, the singular set $\Sigma(F)$ is modified by isotopy,  as well as the above listed moves.
\end{theorem}

\begin{proof}  Let $F\co M\times [0,1]\to \R^2\times [0,1]$ be a generic homotopy, and let $\pi\co M\times [0,1]\to [0,1]$ denote the projection to the second factor. If $\pi|_{A_I(F)}$ does not have critical points on the level $M\times \{t_0\}$, then for sufficiently small $\varepsilon>0$, the singular set $A_I(f_t)$,  parametrized by $t\in [t_0-\varepsilon, t_0+\varepsilon]$,  is modified by an ambient isotopy.  Thus, it remains to study modifications of the singular set of $f_t$ corresponding to critical points of the generic functions $\pi|_{A_I(F)}$.  We claim that $\pi|_{A_I(F)}$ has no critical points when $I=\{1\}$. 

\begin{lemma}\label{le:4}  The map $\pi|_{A_1(F)}$ is a submersion.
\end{lemma}
\begin{proof}   
Over the set $A_1(F)$ of critical points, there is a well-defined kernel bundle $K_1(F)$ of $dF$. In fact, over $A_1(F)$ there is a splitting  
\[
       T(M\times [0,1])|_{A_1(F)} \cong K_1(F)|_{A_1(F)} \oplus TA_1(F). 
\]
Assume that there is a critical point $p\in A_1(F)$ of the function $\pi|_{A_1(F)}$. Then $T_p(A_1(F))$ is in the kernel of $d_p\pi$. On the other hand, the projection $d_p\pi$ coincides with the composition 
\[
      T_p(M\times [0,1])\longrightarrow T_{F(p)}(\R^2\times [0,1])\longrightarrow T_{\pi(p)}([0,1])
\]
of $d_pF$ and the differential of the projection $\R^2\times [0,1]\to [0,1]$ onto the second factor. Since $K_1(F)|_p$ is in the kernel of $d_pF$, it follows that $K_1(F)|_p$ is in the kernel of $d_p\pi$. To summarize, we have shown that $T_p(M\times [0,1])$ is in the kernel of $d_p\pi$, which contradicts the fact that $\pi$ is a submersion.  
\end{proof}

\noindent Let us now consider critical points of the function $\pi|_{A_2(F)}$. 

\begin{lemma}\label{le:5} Let $p\in A_2(F)$ be a critical point of $\pi|_{A_2(F)}$. Then $p$ is a critical point of $\pi|_{\Sigma(F)}$. 
\end{lemma}
\begin{proof}  As above, over $A_2(F)$, there is a well-defined kernel bundle $K_1(F)$. 
\iffalse
which fit the exact sequence
 \[
      0\longrightarrow K_1(F) \longrightarrow T(M\times [0,1])|_{A_2(F)}\longrightarrow F^*T(\R^2\times [0,1])|_{A_2(F)}
      \longrightarrow Q_1(F) \longrightarrow 0. 
\]
\fi
Let $L$ denote the vector subbundle of $T(M\times [0,1])|_{A_2(F)}$ given by $K_1(F)\cap T(\Sigma(F))$. It follows that $\dim L=1$, and there is a splitting
\[
     T(\Sigma(F))|_{A_2(F)} \cong L \oplus T(A_2(F)). 
\]
Since $L_p$ belongs to the kernel $K_1(F)|_p$ of $d_p\pi$, it belongs to the kernel of $d_p\pi|_{\Sigma(F)}$.
On the other hand, if $p$ is a critical point of $\pi|_{A_2(F)}$, then $T_p(A_2(F))$ is also in the kernel of $d_p\pi|_{\Sigma(F)}$. Thus, the point $p$ is a critical point of $\pi|_{\Sigma(F)}$. 
\end{proof}

By Lemma~\ref{le:5}, if $p$ is a critical point of $\pi|_{A_2(F)}$, then $p$ is also a critical point of the function $\pi|_{\Sigma(F)}$. If the index of the critical point $p$ is $0$, then $p$ corresponds to the appearance (birth) of a wrinkle singularity in $\Sigma(f_t)$. A critical point of index $1$ corresponds to the cusp merge move or its inverse, while a critical point of index $2$ corresponds to the disappearance (death) of a wrinkle singularity. 

The critical points of $\pi$ restricted to the submanifold of double points of $A_{11}(F)$ correspond to Reidemeister-II fold crossings. All points of $A_{12}(F), A_{111}(F)$,  and $A_{3}(F)$ are critical in the sense that the differential of $\pi|_{A_I(F)}$ in these cases vanishes. 
It remains to observe that points of $A_{12}(F)$ correspond to cusp-fold crossings,  $A_{111}(F)$ correspond to  Reidemeister-III fold crossings, and $A_3(F)$ correspond to swallowtail singularities.  
\end{proof}

\end{document}
